\section{Introduction}
\label{sec:introduction}

Conditional inference separates feature selection from threshold optimization to
reduce split selection bias. Classical
CART trees choose both a feature and a threshold from the same
label-dependent search, which produces the documented preference for
features with many admissible cutpoints
\citep{Breiman1984ClassificationAndRegressionTrees,LohShih1997SplitSelectionMethods,KimLoh2001ClassificationTreesUnbiasedMultiwaySplits,Loh2002RegressionTreesUnbiasedVariableSelectionInteractionDetection,Shih2004NoteOnSplitSelectionBiasClassificationTrees,StroblBoulesteixAugustin2007UnbiasedSplitSelectionGini,Strobl2007BiasRFVariableImportance}.
In the conditional inference framework \citep{Hothorn2006UnbiasedRecursivePartitioning,HothornHornikVandeWielZeileis2006LegoSystemConditionalInference},
Stage~A tests features for association with the response at the
current node and selects among them; Stage~B then optimizes a threshold within
the selected feature. That separation targets split selection bias, but
conditional inference can require large permutation budgets and exhaustive
threshold scans at every node.
The same conditional inference framework also motivated later work on
conditional variable importance \citep{Strobl2008ConditionalVIM}.
A related line of work uses forests directly for variable screening and ranking,
rather than only for final prediction
\citep{GenuerPoggiTuleauMalot2010VariableSelectionUsingRandomForests}.

This paper evaluates conditional inference trees (CIT) and conditional
inference forests (CIF), as implemented in the software of
\citet{MilletichDownesHirst2026citrees}, as feature rankers. It asks how their
rankings compare with established rankers, and what the Bonferroni corrections
inside their split tests, the feature correction in the Stage~A feature test
and the threshold correction in the Stage~B threshold test, cost and change.
For each seed and cross-validation fold, a method produces a ranking from the
training fold. We then fit fixed downstream learners on top-$k$ prefixes of that
ranking, score them on the held-out fold, and summarize the scores across
datasets. We also ask whether the runtime hyperparameters of CIT and CIF reduce
computation while preserving ranking quality, and how well CIF recovers known
informative features in high-dimensional and synthetic designs. This
feature-selection benchmark is separate from work on
trees and forests as direct
predictors or causal estimators
\citep{Breiman2001RandomForests,WagerAthey2018CausalForests,AtheyTibshiraniWager2019GeneralizedRandomForests}.
Because several methods have tunable configurations, the benchmark chooses one
configuration per method and task. As a sensitivity check, leave-one-dataset-out
(LODO) analyses repeat the configuration choice after holding out each dataset, on
the complete-case and all-method panels defined in \cref{sec:experiments-analysis-plan}.

\paragraph{Contributions.}
\begin{enumerate}
\item A benchmark of CIT and CIF against 15 classification and 16 regression
rankers on 31 real datasets with five seeds, scored by fixed downstream learners
on top-$k$ prefixes (\cref{sec:experiments-real}). CIF is the strongest
conditional inference ranker in the study, ahead of cforest, ctree, and single
CIT, with bootstrap intervals for the paired score difference that exclude zero
in classification. The comparators include random forests (RF) and recursive
feature elimination with random forests (RF-RFE), and positions are on the
all-method panels unless stated. CIF ranks tied 3rd of 17 in classification,
behind RF-RFE and XGBoost and level with RF; 2nd of 18 in regression, behind
ExtraTrees. We read the 0.05 gap in mean rank to XGBoost as a tie. On the
complete-case panels CIF is tied 6th of 17 in classification, behind all three boosting families, and tied
5th of 18 in regression, where the Friedman test gives $p=0.072$ and the
ordering is descriptive.
\item The negative results of that comparison (\cref{sec:boundary,sec:synthetic-recovery,sec:mechanism}). On the
synthetic designs, where the informative features are known, CIF sits in the
lower half on recovery, 8th of 15 in classification and 11th of 16 in
regression. It is weakest at small $k$, ranking 6th to 8th at $k \le 10$ on the
complete-case panels against 3rd or 4th at $k=100$
(\cref{fig:k-trajectory,fig:regression-k-trajectory}). A random forest ranked by
out-of-bag permutation importance does no better than CIF or the
impurity-ranked forest (\cref{tab:rf-oobperm-benchmark}).
\item An account of the cost of the corrections
(\cref{sec:results-practical,sec:results-stageB}). Under the minimum budget,
which gives each of the $m_t$ features at a node $\lceil m_t/\alpha\rceil$
permutations and each of its $K$ candidate thresholds $\lceil K/\alpha\rceil$,
the Stage~A feature test costs, when every test runs, of order $m_t^2/\alpha$
permutations per node and the per-threshold rule of Stage~B of order
$K^2/\alpha$. The runtime ablations
place most of the fitting time in the threshold search; the feature term
dominates on the widest data, where a single tree on the 5,000 features of
gisette takes 29 minutes.
\item A max-type rule for the Stage~B threshold test, an opt-in alternative to
the default per-threshold rule (\cref{cor:stageB-maxt,sec:results-stageB}). At
a fixed budget of 999 permutations the per-threshold rule cannot reject once a
node has more than about 50 candidates. The max-type rule reaches the nominal
level with about $999K$ statistic evaluations, has realized size closer to
nominal, loses about 0.006 in power when the two rules are compared at equal
realized size, and reduces the preference of the per-threshold rule for
features with few distinct values.
\item As a secondary result, the null size of the adaptive stopping rule, exact
under A0.1 to A0.3 at fixed $B$ for a continuous null
(\cref{prop:adaptive-size,app:assumptions}). The rule is inert at the minimum budget the
benchmark uses, because that budget equals the smallest budget at which the
rule may stop, and it acts only at the larger auto budget, the library default.
\end{enumerate}

The theory establishes the validity of each split test at a fixed node,
conditional on the node sample; the cardinality behavior of the rankings is
established empirically. Under the minimum budget every rejecting candidate
returns the smallest attainable $p$-value, so the benchmarked procedure selects
features and thresholds by the order of the observed statistics behind the
feature and threshold corrections (\cref{sec:method}); the guarantees rest on
the standard fixed-node permutation result, stated as a cited lemma
(\cref{lem:plusone-superuniform}), and the benchmark measures what that
procedure delivers in practice. \Cref{sec:method} describes the tree-growing
method, \cref{sec:theory} the fixed-node results, \cref{sec:experiments} the
experimental design, \cref{sec:results,sec:boundary} the results, and
\cref{sec:discussion,sec:conclusion} their interpretation and limits.
